% ===== §2 The Phenomenology and Diagnostics of the Everyday =====
\section{The Phenomenology and Diagnostics of the Everyday}
\label{sec:diagnostics}

\noindent
Before a theory of everyday aesthetics can propose anything, it owes two things that are usually run together and are better kept apart: a description of how the everyday actually shows itself to perception, and a diagnosis of what has gone wrong with it and what a healthy relation to it would be. The first is phenomenological and withholds judgement; the second is diagnostic and does not. This section supplies both in that order, because the proposals of the later parts of the paper are answerable to a diagnosis, and a diagnosis that has not first looked carefully at the phenomenon it means to treat is a prejudice with clinical pretensions. What the description will show is an aesthetic life of low amplitude and high frequency, pervasive and endangered in the same breath. What the diagnosis will add is an account of the danger, of the ground on which the everyday nonetheless remains the decisive field, and of what it would be for the dynamics of a life and a relation to be well rather than ill within it.

\subsection{The phenomenology of the everyday}
\label{sub:diag-phenomenology}

The everyday shows itself, first, in a register beneath the threshold at which experience ordinarily announces itself as aesthetic. It does not arrive framed. The quality of morning light on a wall, the particular sound a familiar door makes, the warmth of the first cup held in both hands, the way a person one lives with sets down a book: these present themselves not as objects offered for contemplation but as the barely-noticed texture of being somewhere, doing something, with someone. Their amplitude is low. No single instance of everyday beauty makes the demand upon attention that a painting in a gallery is institutionally arranged to make; each could be, and mostly is, passed over without remark. Their frequency, correspondingly, is enormous. Where the encounter with art is rare and set apart, the everyday delivers its small aesthetic events continuously, through every waking hour, so that the total aesthetic weight of a life is carried far more by this constant low murmur than by the few occasions on which one stands deliberately before something made to be seen.

The everyday shows itself, second, as transient in a way that is constitutive rather than incidental. The light that is precisely this on the wall now will not be this again; even were the same people to gather tomorrow in the same room with the same things, the gathering would not be the same gathering, for the light would have moved, and the people would have altered, and something would have shifted that cannot be shifted back. The tradition of tea gave this its exact name. \zh{一期一会}, one time, one meeting: a phrase drawn from the practice of tea and crystallised by Ii Naosuke, marking each gathering as singular and unrepeatable, and doing so, as its commentators are careful to note, in a register that is relational before it is contemplative, concerned with the shared event between persons and not merely with the inner state of one \citep{saito2007}. The everyday is the standing home of this transience. Its beauty is very largely the beauty of what is passing, and this is not a melancholy attaching to the everyday from outside but a feature of how the everyday is given at all.

The everyday shows itself, third and most consequentially, under a constant tendency to recede from view. What is encountered continuously is precisely what perception learns to stop registering. This is not a failure that befalls some perceivers and spares others; it is the ordinary working of a perceptual system that economises by ceasing to attend to what it has learned to predict. The phenomenological description must therefore include, as one of the everyday's own features, its liability to become invisible to the very perceiver it surrounds. It is here that description shades into diagnosis, for this liability is at once a fact about how the everyday appears and the opening through which its characteristic pathologies enter.

\subsection{Three pathologies}
\label{sub:diag-pathologies}

The everyday fails perception in three linked ways, and naming them precisely matters, because the models proposed later are, among other things, answers to these three.

The first pathology is \emph{habituation}. The continuous and the familiar are what perception ceases to see. Russian Formalism gave this its sharpest early statement: everyday life, Shklovsky wrote, produces automatisation, a condition in which perception goes mechanical, and habitualisation \emph{devours} the objects of a life, its work and its clothing and its furniture and even, in his unsparing list, one's spouse \citep{shklovsky1965}. The observation is exactly right and the paper takes it as the first term of its diagnosis. But it is worth marking at once where this paper will part from the remedy Formalism attached to the observation, since the divergence sets the whole direction of what follows. Shklovsky's cure was defamiliarisation: the work of art makes the familiar strange, interrupts the automatism, and restores the freshness of a perception that ordinary life had dulled. The cure locates the aesthetic power in a technique applied to the object, and it locates it, tellingly, in \emph{art}, so that the everyday is redeemed only by being made, through artistic device, to stop being everyday. The path taken here is the contrary one. The remedy for habituation will not be to make the ordinary strange but to cultivate a relational capacity by which the ordinary is generated anew in its very ordinariness. One does not defeat the deadening of the light on the wall by rendering the light bizarre; one defeats it by the generative work of a perception, and above all a shared perception, that meets the light as this light, now, once. Shklovsky diagnosed the disease with precision and prescribed an escape from the everyday; this paper prescribes a deeper entry into it.

The second pathology is \emph{numbing}, which is habituation's affective completion. Where habituation is the failure to perceive, numbing is the failure to be moved by what is perceived, the flattening of affective response until the small aesthetic events of the day register, if at all, without consequence. Numbing is the more dangerous because it can wear the costume of depth: a perceiver gone affectively flat may mistake the flatness for equanimity or maturity, and a theory of the everyday must be able to tell the deadening of response apart from any genuine composure. The mark of the difference is generative: composure remains open to being moved and simply holds the movement well, whereas numbing has closed the channel by which the everyday could move one at all.

The third pathology is \emph{symbolic decoration}, and it is the subtlest, because it looks like the cure rather than the disease. Faced with an everyday gone dead, a life may reach for the symbol: it may photograph the meal, narrate the moment, arrange the ordinary into a legible sign of itself to be displayed or archived. The gesture promises to rescue the everyday from invisibility by marking it as significant. What it frequently accomplishes instead is the substitution of the sign for the perception, so that one possesses the token of having-seen in place of the seeing. This third pathology cannot be properly analysed here, because its analysis requires the theory of the symbolic order and its two fates that \xref{sec:symbolic} will supply; the symbol's intervention in perception is not in itself a pathology, and only a certain fate of that intervention is. It is entered here as the third term of the diagnosis and left, deliberately, as a debt to be paid when the resources to pay it are in hand.

\subsection{The everyday as the field of contingency}
\label{sub:diag-contingency}

Diagnosis so far has been negative, a naming of failures. It must now become positive, and it becomes positive by saying what the everyday is such that these are its failures. The everyday is the field of contingency. This is not a fresh coinage but the specific character taken, at the scale of daily life, by the category of the field that earlier papers of this series established, the conditional structure within which relational production occurs. At this scale the field's defining character is that what happens in it is not arranged in advance. The everyday is precisely the region of a life that is not staged, not brought to a ceremonial peak, not curated into the frame of art or ritual; and its contingency, far from being a deficiency it should overcome by aspiring to the condition of the framed and the arranged, is the very condition of its generativity. A field in which everything is arranged admits only execution: what is fully staged can be performed, and performed well or badly, but it cannot generate, because generation requires that something occur which was not laid down in advance. The everyday's refusal of arrangement is therefore not the absence of aesthetic order but the presence of the room in which aesthetic order can be generated rather than merely enacted. One time, one meeting names this too: the singular and unrepeatable gathering is aesthetically decisive exactly because its precise configuration was contingent, never to recur, and so is a site of generation and not of reproduction. The pathologies of the previous section can now be understood in their unity. Each is a way in which the field of contingency is prevented from generating: habituation forecloses the contingent by predicting it away, numbing forecloses it by rendering it inconsequential, and decorative symbolisation forecloses it by settling it prematurely into a sign. The disease, in each of its three forms, is the arrest of a generativity that the contingency of the everyday was the standing opportunity for.

\subsection{The dynamics of flourishing}
\label{sub:diag-flourishing}

The positive diagnosis can be completed by saying what it would be for a life and a relation to be well within the field of contingency, and here this paper draws on the account of happiness given earlier in the series and does not rebuild it. That account declined to treat happiness as a state to be reached and treated it instead as a quality of a dynamical process, a matter of how the coupled systems of a person and a relation move through time rather than of any terminus at which they might arrive \citep{paperxix}. Set in the present frame, the claim is that a healthy everyday is one in which the dynamics of an individual system and of a relational system develop, within the field of contingency, in a direction that tends toward flourishing; and that the aesthetic life of the everyday is the sensuous medium in which this development is largely carried. The small continuous aesthetic events of the day are not ornaments upon a life whose real business lies elsewhere. They are among the principal couplings through which a person and a relation move each other, the medium in which the trajectory toward or away from flourishing is actually traced. This is why the arrest named by the three pathologies is not a merely aesthetic loss. When the field of contingency is prevented from generating, the medium through which the coupled dynamics of a life could tend toward flourishing is deadened, and the persons within it are left to move through a day that has stopped moving them. The stakes of an aesthetics of the everyday are, in the end, the stakes of this dynamics: to keep the ordinary generative is to keep open the medium in which two people, turned together toward their shared and unrepeatable days, may go on becoming well.
